\documentclass[
  spanish,
  letterpaper, oneside
]{article}

\def\documenttitle {Bienvenida, introducción y objetivo de Derecho a la Seguridad Social}
\def\documentsubtitle {Encuadre inicial, Derecho Social y ruta operativa de protección}
\def\documentsubject {Licenciatura en Derecho}

\def\documentauthor {Martin Jonathan de la Cruz}
\def\coursename {Derecho a la Seguridad Social}
\def\coursecode {025}

\def\universityname {Universidad Abierta y a Distancia de México}
\def\universityfaculty {Licenciatura en Derecho}
\def\universitydepartment {Derecho a la Seguridad Social}
\def\universitydepartmentimage {departamentos/UnADM}
\def\universitydepartmentimagecfg {height=1.57cm}
\def\universitylocation {Roma Norte, Ciudad de México}

\def\authortable {
  \begin{tabular}{ll}
    \textbf{Alumno:}
    & \begin{tabular}[t]{l}
      Martin Jonathan de la Cruz \\
    \end{tabular} \\
    Matricula: & ES2611202040 \\
    Figura docente: & Cristian Ortega Barrera \\
    Semestre/Bloque: & 2 / 1 \\
    & \\
    \multicolumn{2}{l}{Fecha de realizacion: \today} \\
    \multicolumn{2}{l}{\universitylocation}
  \end{tabular}
}

\input{template}
\usepackage{helvet}
\renewcommand{\familydefault}{\sfdefault}
\usepackage{array}
\usepackage{longtable}
\usepackage{pdflscape}
\usepackage{tikz}
\usetikzlibrary{arrows.meta,positioning,calc,matrix,fit,shapes.geometric,shadows.blur}
\setcitestyle{authoryear,open={(},close={)}}

\def\coverwatermarkenabled {true}
\def\coverwatermarkimage {img/departamentos/UnADM.pdf}
\def\coverwatermarkopacity {0.16}
\def\coverwatermarkwidth {0.72\paperwidth}
\def\coverwatermarkxshift {0cm}
\def\coverwatermarkyshift {-0.35cm}

\newcommand{\insertcoverwatermark}{%
  \ifthenelse{\equal{\coverwatermarkenabled}{true}}{%
    \AddToShipoutPictureBG*{%
      \begin{tikzpicture}[remember picture,overlay]
        \node[
          opacity=\coverwatermarkopacity,
          inner sep=0pt,
          xshift=\coverwatermarkxshift,
          yshift=\coverwatermarkyshift
        ] at (current page.center) {%
          \includegraphics[width=\coverwatermarkwidth]{\coverwatermarkimage}%
        };
      \end{tikzpicture}%
    }%
  }{}%
}

\begin{document}

\insertcoverwatermark
\templatePortrait
\templatePagecfg
\onehalfspacing

\begin{abstractd}
  Este documento integra la fuente local consolidada de bienvenida, introducción
  y objetivo de la asignatura de \textit{Derecho a la Seguridad Social}. Su
  propósito es conservar, en formato académico y verificable, el sentido de la
  materia como rama del Derecho Social, el encuadre docente, los conceptos clave
  y una ruta operativa para analizar contingencias, instituciones, prestaciones y
  vías de defensa de personas derechohabientes.
\end{abstractd}

\templateIndex
\templateFinalcfg

\section{Bienvenida y sentido de la materia}

La materia \textit{Derecho a la Seguridad Social} se orienta a formar competencias para proteger derechos frente a imprevistos como enfermedades, accidentes, vejez, cesantía, desempleo o invalidez. Estas eventualidades reducen los recursos de las personas justo cuando necesitan apoyo para conservar condiciones dignas de vida.

El enfoque central es que la seguridad social no debe verse como ayuda discrecional, sino como un sistema jurídico de protección. Su estudio permite comprender cómo se crean derechos, obligaciones e instituciones para que las personas accedan a servicios o recursos que normalmente no podrían pagar por sí mismas.

\section{Presentación del docente}

El docente se presenta como \textbf{Cristian Ortega Barrera}, Maestro y Licenciado en Derecho por la Facultad de Derecho de la Universidad Nacional Autónoma de México, con aproximadamente 23 años de experiencia profesional en el sector público y privado. Este dato proviene de la nota local de bienvenida y se conserva como trazabilidad interna de la materia; no se usa como fuente bibliográfica final porque no sustituye fuentes académicas recuperables.

\section{Ubicación en el Derecho Social}

El Derecho a la Seguridad Social pertenece al \textit{Derecho Social}, junto con ramas como el Derecho Laboral y el Derecho Agrario. Esta ubicación es relevante porque el Derecho Social no se limita a regular relaciones entre partes formalmente iguales; atiende desigualdades materiales, situaciones de vulnerabilidad y necesidades de protección frente a riesgos de la vida.

La seguridad social cumple una función redistributiva: equilibra, al menos parcialmente, la relación entre personas con riqueza o estabilidad económica y personas que enfrentan enfermedad, desempleo, vejez, accidentes o incapacidad. Por ello, su estudio exige vincular norma, institución, contingencia y prestación.

\section{Introducción a la unidad didáctica}

En México, la seguridad social se vincula históricamente con luchas sociales como la Revolución Mexicana de 1910 y con la construcción del Derecho Laboral, Agrario y Social. Aunque inició con formas asistencialistas y después se consolidó como protección vinculada al empleo formal, las crisis sanitarias, económicas y laborales obligan a pensarla como un mecanismo de protección más amplio.

El derecho a la seguridad social puede analizarse como derecho humano y como obligación jurídica para ciertos sujetos, especialmente patrones. La tensión contemporánea consiste en equilibrar la lógica contributiva del empleo formal con la necesidad de proteger a personas fuera de ese empleo, con trayectorias discontinuas o en situación de vulnerabilidad.

\section{Objetivo de la asignatura}

El objetivo de la materia es identificar y comprender los aspectos teóricos y técnicos de los derechos de la Seguridad Social y de aquellos que emanen de este, para su debida aplicación. En términos operativos, esto implica saber detectar una contingencia, ubicar el régimen aplicable, identificar institución competente, verificar requisitos y elegir una vía de defensa eficaz.

\section{Conceptos clave de inicio}

\begin{longtable}{>{\bfseries}p{0.30\textwidth}p{0.62\textwidth}}
\toprule
Concepto & Sentido dentro de la asignatura \\
\midrule
Derecho a la Seguridad Social & Derecho humano y disciplina jurídica orientada a proteger frente a contingencias que afectan salud, ingreso, subsistencia o retiro. \\
Derecho Social & Rama que atiende relaciones asimétricas y busca proteger a personas o grupos en vulnerabilidad. \\
Personas derechohabientes & Personas que pueden acceder a servicios o recursos derivados de sistemas de seguridad social. \\
Contingencia & Evento como enfermedad, accidente, vejez, desempleo o invalidez que activa necesidades de protección. \\
IMSS e ISSSTE & Sistemas institucionales mexicanos que forman parte del estudio y defensa de la seguridad social. \\
Ahorro para el retiro & Mecanismo relacionado con protección de vejez, retiro digno y continuidad de ingresos. \\
\bottomrule
\end{longtable}

\section{Ruta inicial de aprendizaje}

Para convertir el derecho reconocido en protección efectiva, cada caso debe ordenarse con una ruta mínima:

\begin{enumerate}
  \item Identificar la contingencia: enfermedad, accidente, vejez, desempleo, invalidez, cesantía u otra afectación.
  \item Determinar régimen e institución competente: IMSS, ISSSTE, sistema de retiro u otro mecanismo aplicable.
  \item Verificar prueba mínima: afiliación, salario, semanas cotizadas, dictamen, beneficiarios, constancias y vigencia de derechos.
  \item Distinguir responsable o vía: institución, patrón, sistema de ahorro, recurso administrativo, vía laboral o protección constitucional.
  \item Priorizar urgencia: salud, ingreso, subsistencia, plazos breves y continuidad de cobertura.
\end{enumerate}

\section{Tensiones de análisis}

La seguridad social debe defenderse como derecho social indispensable, pero también evaluarse con criterios de eficacia real. No basta con que una norma reconozca una prestación si existen barreras de afiliación, trámites excesivos, falta de prueba, errores de cotización, baja calidad institucional o ausencia de financiamiento sostenible.

Por ello, la materia debe trabajar una tensión permanente: ampliar la protección sin ignorar sostenibilidad, cobertura efectiva, calidad del servicio y capacidad administrativa. El enfoque equilibrado forma juristas capaces de defender personas derechohabientes y, al mismo tiempo, diagnosticar fallas sistémicas.

\section{Términos relacionados}

\begin{itemize}
  \item Derecho Social.
  \item Derecho Laboral.
  \item Derecho Agrario.
  \item IMSS e ISSSTE.
  \item Ahorro para el retiro.
  \item Derechohabientes.
  \item Cesantía, vejez, desempleo e invalidez.
  \item Cobertura efectiva y sostenibilidad.
\end{itemize}

\section{Lectura editorial}

La unificación editorial con \textit{Garantías Constitucionales} no traslada contenido temático exclusivo entre materias. Solo adopta el patrón formal y didáctico: diagnóstico inicial, distinción conceptual, ruta de análisis, validación con fuentes verificadas y cierre con postura profesional. En Seguridad Social, la capa local dominante sigue siendo el análisis de contingencias, instituciones, prestaciones y defensa de personas derechohabientes.

\section{Conclusión}

El \textit{Derecho a la Seguridad Social} convierte riesgos previsibles de la vida en responsabilidades jurídicas y sociales compartidas. Su valor no se agota en estudiar instituciones como IMSS, ISSSTE o sistemas de retiro; consiste en comprender cómo esas instituciones pueden proteger efectivamente salud, ingreso, vejez, subsistencia y dignidad. Mi postura es que la materia debe abordarse como derecho humano, sistema institucional y herramienta práctica de defensa, siempre con verificación normativa, probatoria y contextual.

\clearpage
\nocite{unadmSitioWeb,unadmMallaDerecho2024,cpeum2026,lss2026,lissste2026,imssSitio2026,isssteSitio2026,consarSitio2026}
\bibliography{derecho-a-la-seguridad-social}

\end{document}
